% ============================================================
%  Capítulo -- Tecnologías utilizadas
% ============================================================

\chapter{Tecnologías utilizadas}

En este capítulo se describen las tecnologías seleccionadas para el desarrollo de
OmniLitter y la justificación de su elección. La selección se ha realizado teniendo
en cuenta tres restricciones principales: el alcance de un Trabajo de Fin de Grado
desarrollado por dos personas, la necesidad de disponer de una aplicación móvil
capaz de trabajar sin conexión y la conveniencia de mantener un stack homogéneo
basado en TypeScript.

\section{Criterios de selección}

Los criterios utilizados para seleccionar las tecnologías han sido los siguientes:

\begin{itemize}
  \item \textbf{Homogeneidad tecnológica:} se prioriza el uso de TypeScript tanto
    en el frontend como en el backend para reducir el cambio de contexto durante
    el desarrollo.
  \item \textbf{Capacidad offline:} la aplicación móvil debe poder crear sesiones,
    observaciones y fotografías sin depender de la conexión a Internet.
  \item \textbf{Madurez del ecosistema:} se eligen tecnologías ampliamente usadas,
    con documentación suficiente y comunidad activa.
  \item \textbf{Facilidad de despliegue:} el sistema debe poder desplegarse en
    servicios accesibles para un proyecto académico, como Render o Microsoft Azure,
    a partir de imágenes de contenedor.
  \item \textbf{Escalabilidad razonable:} aunque el proyecto no persigue una escala
    industrial, el diseño debe permitir incorporar nuevos grupos, campañas y datos
    de campo sin rehacer la arquitectura.
\end{itemize}

\paragraph{Tecnologías empleadas en los prototipos}

Durante la fase inicial se construyeron prototipos funcionales con datos en
memoria: un portal web y una simulación de la aplicación móvil en navegador. Estos
prototipos permitieron validar flujos de usuario con el cliente antes de trasladar
las decisiones principales a la implementación final.

\begin{table}[H]
  \centering
  \caption{Tecnologías empleadas en los prototipos funcionales.}
  \label{tab:stack_mockups}
  \begin{tabularx}{\textwidth}{lX}
    \toprule
    \textbf{Tecnología} & \textbf{Uso en el proyecto} \\
    \midrule
    React 18 & Construcción de interfaces de usuario para el portal web y el prototipo móvil. \\
    TypeScript & Tipado estático de componentes, rutas, estado y datos de ejemplo. \\
    Vite & Servidor de desarrollo y empaquetado rápido de los prototipos. \\
    Tailwind CSS & Sistema de estilos utilitario para construir pantallas responsive. \\
    React Router & Definición de rutas públicas, rutas protegidas y navegación entre pantallas. \\
    Recharts & Gráficas y visualizaciones estadísticas del dashboard del portal web. \\
    Lucide React / Material Icons & Iconografía de la interfaz. \\
    \bottomrule
  \end{tabularx}
\end{table}

\section{Portal web}

El portal web se plantea como una aplicación \textit{Single Page Application}
desarrollada con React, TypeScript y Vite. Su responsabilidad principal es permitir
a investigadores y administradores consultar observaciones, gestionar usuarios y
grupos, visualizar mapas y exportar datos.

\begin{table}[H]
  \centering
  \caption{Stack seleccionado para el portal web.}
  \label{tab:stack_web}
  \begin{tabularx}{\textwidth}{lX}
    \toprule
    \textbf{Tecnología} & \textbf{Responsabilidad} \\
    \midrule
    React 18 & Construcción de componentes reutilizables y pantallas del portal. \\
    TypeScript & Definición de tipos de dominio, respuestas de API y props de componentes. \\
    Vite & Compilación, servidor de desarrollo y generación del bundle final. \\
    Tailwind CSS & Sistema visual responsive y consistente con los prototipos. \\
    Enrutado propio (History API) & Gestión de rutas públicas y protegidas en el cliente, sin librería de enrutado externa. \\
    Fetch API & Comunicación con la API REST del backend. \\
    Leaflet y React Leaflet & Mapa interactivo de observaciones con agrupación de marcadores (\textit{clustering}) y capa de densidad (\textit{heatmap}). \\
    Recharts & Gráficos del dashboard y de la analítica. \\
    i18next / react-i18next & Interfaz en español e inglés. \\
    \bottomrule
  \end{tabularx}
\end{table}

\section{Aplicación móvil}

La aplicación móvil se plantea con React Native y Expo, utilizando SQLite mediante
el paquete \texttt{expo-sqlite} como base de datos local. El objetivo es que el
usuario pueda completar una jornada de muestreo aunque el dispositivo no tenga
conexión a Internet.

\begin{table}[H]
  \centering
  \caption{Stack seleccionado para la aplicación móvil.}
  \label{tab:stack_mobile}
  \begin{tabularx}{\textwidth}{lX}
    \toprule
    \textbf{Tecnología} & \textbf{Responsabilidad} \\
    \midrule
    React Native & Desarrollo multiplataforma para Android e iOS. \\
    Expo (\textit{bare workflow}, \texttt{expo-dev-client}) & Toolchain de desarrollo y compilación con acceso a módulos nativos. \\
    TypeScript & Tipado de pantallas, servicios, modelos locales y respuestas de API. \\
    \texttt{expo-sqlite} / SQLite & Persistencia local de sesiones, observaciones, residuos, taxonomía y cola de sincronización. \\
    \texttt{expo-location} & Captura de coordenadas GPS asociadas a sesiones y observaciones. \\
    \texttt{expo-camera} / \texttt{expo-image-picker} & Captura de fotografías con la cámara y selección desde la galería. \\
    \texttt{expo-file-system} & Almacenamiento local de las fotografías en el dispositivo. \\
    \texttt{expo-secure-store} & Guardado seguro de credenciales y \textit{token} para el inicio de sesión sin conexión. \\
    \texttt{@react-native-community/netinfo} & Detección del estado de conectividad para disparar la sincronización al recuperar cobertura. \\
    \texttt{react-native-webview} & Mapa de sesión embebido para mostrar el recorrido y el área muestreada. \\
    i18next / react-i18next & Interfaz en español e inglés. \\
    Navegación propia (estado de React) & Cambio entre las pestañas principales y los flujos de captura, sin librería de navegación externa. \\
    \bottomrule
  \end{tabularx}
\end{table}

\section{Backend y persistencia}

El backend centraliza la lógica de negocio, la autenticación, el control de acceso
por grupo y la sincronización de datos procedentes de dispositivos móviles. Se
plantea como una API REST desarrollada con Node.js, Express y TypeScript.

\begin{table}[H]
  \centering
  \caption{Stack seleccionado para backend y persistencia.}
  \label{tab:stack_backend}
  \begin{tabularx}{\textwidth}{lX}
    \toprule
    \textbf{Tecnología} & \textbf{Responsabilidad} \\
    \midrule
    Node.js 20 & Entorno de ejecución del backend. \\
    Express & Definición de rutas, middlewares y controladores REST. \\
    TypeScript & Tipado de servicios, entidades y errores. \\
    PostgreSQL 16 & Base de datos relacional principal del sistema. \\
    Prisma & ORM para modelar entidades, generar migraciones y acceder a los datos. \\
    Zod & Validación de esquemas de entrada, compartida entre cliente y servidor. \\
    JWT (\texttt{jsonwebtoken}) & Autenticación mediante \textit{tokens} de acceso y refresco. \\
    \texttt{bcryptjs} & Almacenamiento de contraseñas mediante \textit{hash}. \\
    Nodemailer & Envío de correo para la recuperación de contraseña. \\
    swagger-ui-express (OpenAPI) & Documentación interactiva de la API en \texttt{/api/docs}. \\
    Helmet, CORS y Morgan & Cabeceras de seguridad, control de origen y registro de peticiones. \\
    Almacenamiento de ficheros & Las fotografías llegan codificadas en \textit{base64} y se guardan en el sistema de archivos con el módulo de ficheros de Node. \\
    \bottomrule
  \end{tabularx}
\end{table}

\section{Despliegue y herramientas auxiliares}

Además de las tecnologías propias de cada componente, el proyecto se apoya en un
conjunto de herramientas de organización, calidad, pruebas y despliegue. La
tabla~\ref{tab:stack_aux} las recoge.

\begin{table}[H]
  \centering
  \caption{Herramientas auxiliares, pruebas y despliegue.}
  \label{tab:stack_aux}
  \begin{tabularx}{\textwidth}{lX}
    \toprule
    \textbf{Tecnología} & \textbf{Responsabilidad} \\
    \midrule
    pnpm \textit{workspaces} & Gestión del monorepo con los paquetes de API, portal web, móvil y código compartido. \\
    TypeScript & Lenguaje común de todos los paquetes, con un paquete compartido de tipos y esquemas. \\
    ESLint y Prettier & Análisis estático y formateo homogéneo del código. \\
    Vitest & Pruebas automatizadas de la API, el portal web y el paquete compartido, con medición de cobertura. \\
    Docker y Docker Compose & Empaquetado y orquestación de PostgreSQL, la API y el portal web. \\
    Nginx & Servidor del portal web dentro de su contenedor. \\
    Git, GitHub y GitHub Actions & Control de versiones e integración continua: validación, publicación de imágenes y generación de versiones. \\
    EAS (Expo Application Services) & Compilación del APK de Android de la aplicación móvil. \\
    Render y Microsoft Azure & Despliegue del sistema en los entornos de preproducción y producción. \\
    \bottomrule
  \end{tabularx}
\end{table}

\section{Alternativas consideradas}

\paragraph{Justificación de SQLite frente a WatermelonDB}

Aunque inicialmente se valoró WatermelonDB como solución local \textit{offline-first},
se ha decidido emplear SQLite directamente para ajustar mejor el alcance del TFG.
SQLite permite cubrir los requisitos principales de la aplicación móvil: almacenar
sesiones, observaciones, fotografías pendientes y una cola de sincronización local.

La decisión reduce la complejidad técnica, evita depender de una capa adicional de
sincronización y facilita que los autores controlen de forma explícita el esquema
local y el flujo de subida de datos. El patrón \textit{offline-first} se mantiene,
pero se implementa mediante una tabla local de cola de sincronización y servicios
propios de lectura/escritura sobre SQLite.
